% ═══════════════════════════════════════════════════════════════════
%  Section 7 — Conclusion
% ═══════════════════════════════════════════════════════════════════

\label{sec:conclusion}

This thesis presented a controlled experimental platform for evaluating LLM bargaining behaviour in alternating-offer bilateral negotiation under private economic constraints.  The platform enforces hard budget and cost limits, provides a configurable deterministic settlement gate, logs every agent action at sufficient granularity to separate LLM-generated decisions from mechanically enforced settlements, and connects micro-level bargaining behaviour to macro-level market outcomes within a single system.  Using this platform, we conducted six experiments totalling 2,770 negotiation sessions with a single 3B-parameter language model (\texttt{llama3.2:3b}).

The experiments produced two categories of findings.  First, LLM agents reproduced several structural predictions of classical bargaining and market theory: deal success was determined by the width of the zone of possible agreement, market prices stabilised rapidly with persistent within-tick dispersion, the buyer-first protocol produced systematic surplus asymmetry, exogenous shocks shifted prices and liquidity in theoretically predicted directions, and surplus-maximising matching improved deal rates and allocative efficiency relative to random matching.  These outcomes arise primarily from the negotiation protocol and private constraints rather than from agent sophistication---a 3B-parameter model, assisted by a feasibility gate that drives approximately 40\% of acceptances, is sufficient to produce qualitatively correct market structure.

Second, agents failed to exhibit two well-documented human bargaining heuristics: anchoring bias and deadline-driven concession acceleration.  These null results are informative but heavily confounded.  The anchoring null may reflect the prompt's provision of explicit surplus information, which removes the uncertainty that drives anchoring in human negotiators.  The deadline null is confounded by the feasibility gate (which settles deals before deadlines approach) and by prompt instructions that discourage rejection.  Neither confound has been ablated; the null results characterise LLM behaviour under these specific experimental conditions, not LLM negotiation capability in general.

A recurring methodological theme is the tension between engineering reliability and experimental cleanliness.  The feasibility gate was necessary to make small-model experiments viable---without it, wide-ZOPA scenarios produce near-zero deal rates due to LLM output quality---but it simultaneously masks the multi-round dynamics that several experiments were designed to study.  More broadly, the prompt functions as the agent's entire cognitive architecture, making all results inherently prompt-conditional.  These observations apply beyond our specific platform: researchers using LLM agents in structured economic settings should treat engineering interventions and prompt design choices as first-order experimental variables, disclose their involvement rates, and conduct ablation studies to separate agent capability from system-level enforcement.

Three directions for future work are most pressing.  First, gate and prompt ablation---disabling the gate and separately removing the accept-condition block, prefer-counter steering, and deadline-salience language---would reveal which null results are artifacts and which reflect intrinsic model behaviour.  Second, repeating key experiments with models of different sizes and architectures would test whether anchoring and deadline effects emerge at larger scales, addressing the most important external-validity question.  Third, the platform supports two additional experiment families---reputation and repeated interaction, and communication strategy ablation---for which infrastructure is implemented but empirical evaluation is deferred to future work.

The platform is a testbed, not a finished product.  Its value lies in enabling systematic, reproducible comparison of LLM bargaining behaviour against well-characterised theoretical predictions---including honest reporting of where that behaviour aligns with theory and where it does not.
